\section{Discussion and Limitations}
\label{sec:discussion}

\paragraph{Simulator-conditional claims}
Following~\cite{stoffregen2020sim2real}, we do not equate synthetic event
performance with real sensor performance. Our detection deltas between
channels are robust \emph{within} the simulator family we study, and we
report sensitivity bands to indicate which conclusions survive parameter
perturbation. When FENCE-class hardware data becomes public, FLIR-STED and
our pipeline are positioned as the ready-made comparison baseline.

\paragraph{Interpolation in the thermal domain}
The failure of flow-guided interpolation (E1) matters beyond our pipeline:
most event simulators inherit visible-light interpolation assumptions. At
low source frame rates and under AGC flicker, brightness constancy is
violated; simple linear blending is safer. Learned interpolators trained on
visible light should be re-audited before thermal use.

\paragraph{Data licensing}
HIT-UAV is CC BY 4.0 and can be redistributed directly. FLIR ADAS v2 is
form-gated; we therefore release conversion code, configurations, and
checksums rather than re-hosting FLIR-derived events, and we publish the
HIT-UAV-derived events directly. KAIST (CC BY-NC-SA) conversions can be
shared under the same licence for academic use.

\paragraph{Realism of the scene model}
Our synthetic scenes are deliberately simple (blobs on drifting textured
backgrounds) so that physics parameters are exactly known for the fidelity
study; they are not a replacement for real thermal texture. The detection
conclusions that rely on realism are therefore anchored on FLIR-STED
(real thermal video), with STED used for controlled parameter studies.

\paragraph{Class collapse}
We keep six FLIR classes and one synthetic object class to keep the study
focused on modality rather than taxonomy; extending to the full FLIR class
set is straightforward.

\section{Conclusion}
\label{sec:conclusion}

Event-based infrared hardware is arriving before any public thermal event
data exists. We have built the first open thermal-to-event conversion
pipeline with thermal-specific AGC de-flickering and bolometer-physics
simulation modes, released the first synthetic and FLIR-derived thermal
event detection benchmarks, and mapped the simulator parameter space with
sensitivity bands. Our three-channel detection study with scene slicing
quantifies where the event channel adds value over thermal frames, and our
interpolation audit warns that visible-light sim-to-real practice does not
transfer unchanged. We position this as a pre-hardware algorithm study:
when event-based infrared cameras ship, the algorithms community should not
start from zero.
